%% main_skeleton.tex -- main file in the house project layout.
%% Copy to the paper folder as the venue-named main file (house practice keeps
%% the kit's stem: iclr2027_conference.tex / neurips_2026.tex / example_paper.tex),
%% next to the venue style files, preamble.tex, references.bib and chapters/.
%% Build:  tectonic -X compile main.tex        (or latexmk -pdf)
%% Default header: ICLR 2027. For NeurIPS / ICML swap the header block using
%% the VENUE NOTES at the top of preamble.tex; keep everything else.

\documentclass{article}
\PassOptionsToPackage{table}{xcolor}      % must precede the venue style
\usepackage{iclr2027_conference,times}
\input{math_commands.tex}                  % official ICLR notation file (\Eqref{})
\newcommand{\VenueLinks}{\hypersetup{colorlinks=true, linkcolor=red,
    filecolor=magenta, urlcolor=cyan, citecolor=cyan}}
%% =====================================================================
%% preamble.tex -- shared preamble in the house style ("Part B").
%% \input this file AFTER the venue header ("Part A", see VENUE NOTES and
%% main_skeleton.tex). It consolidates the packages, macros, colours and boxes
%% the house style uses at all three venues; rationale in
%% references/figures_tables.md and references/math.md.
%% Evidence labels in the comments ([all], [most], [some], [current], [older],
%% [venue: X], [fix], [skill choice]) are explained in SKILL.md.
%% Tested: tectonic -X compile (XeTeX) with iclr2027_conference.sty, with
%% neurips_2026.sty ([main]) and with icml2026.sty (main_skeleton_icml.tex).
%% =====================================================================
%%
%% VENUE NOTES -- keep exactly ONE header in the main file.
%%
%% (A1) ICLR 2027  [default]
%%   \documentclass{article}
%%   \PassOptionsToPackage{table}{xcolor}   % BEFORE the style: eso-pic loads xcolor
%%   \usepackage{iclr2027_conference,times}
%%   \input{math_commands.tex}              % official file; it redefines \eqref,
%%                                          % so write \Eqref{eq:x}, never "Equation \eqref"
%%   \newcommand{\VenueLinks}{\hypersetup{colorlinks=true, linkcolor=red,
%%       filecolor=magenta, urlcolor=cyan, citecolor=cyan}}      % house colours
%%   % \iclrfinalcopy                       % camera-ready only
%%   Cite with \citep{} only ([venue: ICLR]; \citet never).
%%   \bibliographystyle{iclr2027_conference}
%%   Limit: 9 pp at submission / 10 at rebuttal and camera-ready (ICLR 2026 and
%%   2027, verified 2026-09-25). Before the references: \section*{AI Use
%%   Statement} (2027: REQUIRED), \section*{Ethics Statement} and
%%   \section*{Reproducibility Statement} (chapters/8_...statement.tex). 2026-style
%%   kits put "Use of LLMs" last in the appendix instead.
%%   Pasting this Part B at the end of math_commands.tex is the usual ICLR move
%%   ([venue: ICLR]); %% =====================================================================
%% preamble.tex -- shared preamble in the house style ("Part B").
%% \input this file AFTER the venue header ("Part A", see VENUE NOTES and
%% main_skeleton.tex). It consolidates the packages, macros, colours and boxes
%% the house style uses at all three venues; rationale in
%% references/figures_tables.md and references/math.md.
%% Evidence labels in the comments ([all], [most], [some], [current], [older],
%% [venue: X], [fix], [skill choice]) are explained in SKILL.md.
%% Tested: tectonic -X compile (XeTeX) with iclr2027_conference.sty, with
%% neurips_2026.sty ([main]) and with icml2026.sty (main_skeleton_icml.tex).
%% =====================================================================
%%
%% VENUE NOTES -- keep exactly ONE header in the main file.
%%
%% (A1) ICLR 2027  [default]
%%   \documentclass{article}
%%   \PassOptionsToPackage{table}{xcolor}   % BEFORE the style: eso-pic loads xcolor
%%   \usepackage{iclr2027_conference,times}
%%   \input{math_commands.tex}              % official file; it redefines \eqref,
%%                                          % so write \Eqref{eq:x}, never "Equation \eqref"
%%   \newcommand{\VenueLinks}{\hypersetup{colorlinks=true, linkcolor=red,
%%       filecolor=magenta, urlcolor=cyan, citecolor=cyan}}      % house colours
%%   % \iclrfinalcopy                       % camera-ready only
%%   Cite with \citep{} only ([venue: ICLR]; \citet never).
%%   \bibliographystyle{iclr2027_conference}
%%   Limit: 9 pp at submission / 10 at rebuttal and camera-ready (ICLR 2026 and
%%   2027, verified 2026-09-25). Before the references: \section*{AI Use
%%   Statement} (2027: REQUIRED), \section*{Ethics Statement} and
%%   \section*{Reproducibility Statement} (chapters/8_...statement.tex). 2026-style
%%   kits put "Use of LLMs" last in the appendix instead.
%%   Pasting this Part B at the end of math_commands.tex is the usual ICLR move
%%   ([venue: ICLR]); %% =====================================================================
%% preamble.tex -- shared preamble in the house style ("Part B").
%% \input this file AFTER the venue header ("Part A", see VENUE NOTES and
%% main_skeleton.tex). It consolidates the packages, macros, colours and boxes
%% the house style uses at all three venues; rationale in
%% references/figures_tables.md and references/math.md.
%% Evidence labels in the comments ([all], [most], [some], [current], [older],
%% [venue: X], [fix], [skill choice]) are explained in SKILL.md.
%% Tested: tectonic -X compile (XeTeX) with iclr2027_conference.sty, with
%% neurips_2026.sty ([main]) and with icml2026.sty (main_skeleton_icml.tex).
%% =====================================================================
%%
%% VENUE NOTES -- keep exactly ONE header in the main file.
%%
%% (A1) ICLR 2027  [default]
%%   \documentclass{article}
%%   \PassOptionsToPackage{table}{xcolor}   % BEFORE the style: eso-pic loads xcolor
%%   \usepackage{iclr2027_conference,times}
%%   \input{math_commands.tex}              % official file; it redefines \eqref,
%%                                          % so write \Eqref{eq:x}, never "Equation \eqref"
%%   \newcommand{\VenueLinks}{\hypersetup{colorlinks=true, linkcolor=red,
%%       filecolor=magenta, urlcolor=cyan, citecolor=cyan}}      % house colours
%%   % \iclrfinalcopy                       % camera-ready only
%%   Cite with \citep{} only ([venue: ICLR]; \citet never).
%%   \bibliographystyle{iclr2027_conference}
%%   Limit: 9 pp at submission / 10 at rebuttal and camera-ready (ICLR 2026 and
%%   2027, verified 2026-09-25). Before the references: \section*{AI Use
%%   Statement} (2027: REQUIRED), \section*{Ethics Statement} and
%%   \section*{Reproducibility Statement} (chapters/8_...statement.tex). 2026-style
%%   kits put "Use of LLMs" last in the appendix instead.
%%   Pasting this Part B at the end of math_commands.tex is the usual ICLR move
%%   ([venue: ICLR]); \input{preamble} after math_commands.tex is equivalent.
%%
%% (A2) NeurIPS 2026
%%   \documentclass{article}
%%   \PassOptionsToPackage{table}{xcolor}
%%   \PassOptionsToPackage{numbers, compress}{natbib}
%%   \usepackage[main]{neurips_2026}        % camera-ready: [main,final]
%%                                          % ([final] alone stops at \@trackname)
%%   \newcommand{\VenueLinks}{\hypersetup{colorlinks, linkcolor=red,
%%       anchorcolor=blue, citecolor=green}}
%%   Cite with \cite{} (numeric). \bibliographystyle{unsrtnat}.
%%   After the appendix: \clearpage \input{checklist}. 9 content pages.
%%
%% (A3) ICML 2026  [full main file: main_skeleton_icml.tex]
%%   \documentclass{article}
%%   \PassOptionsToPackage{table}{xcolor}   % prevents "Option clash for package
%%                                          % xcolor" when [table]{xcolor} follows the sty
%%   \usepackage[T1]{fontenc}               % B1 loads it too; without T1, XeTeX
%%                                          % silently falls back to Latin Modern
%%   \usepackage{microtype}\usepackage{graphicx}\usepackage{subcaption}
%%   \usepackage{booktabs}\usepackage{hyperref}
%%   \newcommand{\theHalgorithm}{\arabic{algorithm}}
%%   \usepackage{icml2026}                  % blind submission (line-number ruler);
%%                                          % [preprint] arXiv; [accepted] camera-ready
%%   \icmltitlerunning{Short Running Title}
%%   After the \twocolumn[...] title block ALWAYS call \printAffiliationsAndNotice{}
%%   (official template); without it icml2026.sty's end-of-document warning fails
%%   with "Unbalanced write command" (verified). Only [accepted] needs CM italic 9pt
%%   (cmti9); if tectonic times out fetching it, use scripts/tectonic_fetch.sh and
%%   -Z search-path=<dir> (references/venues.md section 2.3).
%%   No \VenueLinks (keep the template's dark-blue links). \cite{} is fine: the sty
%%   maps it to \citep{} (parenthetical, ONE argument; \citep[p.~3]{x} for pages).
%%   8 pp at submission (+1 camera-ready); \section*{Impact Statement} (required)
%%   after \clearpage, before the references. Plots use \columnwidth. The kit is
%%   not bundled: templates/venue/icml/README.md.
%% =====================================================================

% ---------- B1. Encoding --------------------------------------------------
% T1 under BOTH engines: the venue styles select Times ("ptm"), which has no
% TU shape; without T1, XeLaTeX falls back to Latin Modern.
\usepackage{iftex}
\ifPDFTeX
  \usepackage[utf8]{inputenc}
\fi
\usepackage[T1]{fontenc}

% ---------- B2. Packages ([all] core + [current] additions) ---------------
\usepackage{xcolor}                 % [table] comes from \PassOptionsToPackage in Part A
\usepackage{url}
\usepackage{booktabs}
\usepackage{amsmath,amssymb,amsfonts}
\usepackage{mathtools}
\usepackage{amsthm}
\usepackage{thmtools}               % provides thm-restate (restatable theorems)
% no bm [current]: write \mathbf{} (\boldsymbol{} from amsmath if needed); with
% bm the NeurIPS 8pt footnotes need cmmib8, which the offline tectonic cache
% lacked (build stopped; verified).
\usepackage{nicefrac}
\usepackage{microtype}
\usepackage{graphicx}
\usepackage{caption}
\usepackage{subcaption}
\usepackage{enumitem}               % always \begin{itemize}[leftmargin=*]
\usepackage{multirow}
\usepackage{tabularx}
\usepackage{makecell}
\usepackage{wrapfig}                % NeurIPS/ICLR only; never in two-column ICML
\usepackage{tcolorbox}
\tcbuselibrary{skins}
\usepackage{tikz}
\usetikzlibrary{calc,shadows}
\ifdefined\algorithmic\else         % ICML's icml2026.sty already loads these
  \usepackage{algorithm}
  \usepackage{algorithmic}
\fi
\ifdefined\hypersetup\else \usepackage{hyperref}\fi
\ifdefined\VenueLinks \VenueLinks \fi
% No cleveref: write "Figure~\ref{}", "Table~\ref{}", "Section~\ref{}",
% "Appendix~\ref{}", "Equation~\ref{}" by hand (house practice types the word
% by hand [all]; the tie only blocks a line break before the number:
% references/latex_source.md section 3). No method-name macro: type the
% name [all].

% ---------- B3. Latin abbreviations ([all]; trailing space is intentional) ---
% Use inside parentheses: (\ie ...), (\eg ...), (\cf Section~\ref{...}).
% Do not type another comma or space after them.
\providecommand{\ie}{\emph{i.e., }}
\providecommand{\eg}{\emph{e.g., }}
\providecommand{\etal}{\emph{et al.}}
\providecommand{\st}{\emph{s.t. }}
\providecommand{\etc}{\emph{etc.}}
\providecommand{\wrt}{\emph{w.r.t.}~}
\providecommand{\cf}{\emph{cf. }}
\providecommand{\aka}{\emph{aka. }}

% ---------- B4. Co-author / status macros ([all]) ---------------------------
% One macro per co-author, named after that person's initials: rename \coA,
% \coB, \coC to each co-author's initials. Coloured while drafting, black for
% submission and camera-ready; the wrappers stay in the text. House practice
% swaps two commented blocks; here one switch does the same.
\newif\ifdraftcolors
\draftcolorsfalse                  % \draftcolorstrue while drafting
\ifdraftcolors
  \newcommand{\coA}[1]{{\color{orange}{#1}}}
  \newcommand{\revision}[1]{{\color{blue}{#1}}}
  \newcommand{\rebuttal}[1]{{\color{blue}{#1}}}
  \newcommand{\draft}[1]{{\color{purple}{#1}}}
  \newcommand{\warning}[1]{{\color{red}{#1}}}
\else
  \newcommand{\coA}[1]{{\color{black}{#1}}}
  \newcommand{\revision}[1]{{\color{black}{#1}}}
  \newcommand{\rebuttal}[1]{{\color{black}{#1}}}
  \newcommand{\draft}[1]{{\color{black}{#1}}}
  \newcommand{\warning}[1]{{\color{black}{#1}}}
\fi
% ---------- B4b. Key concepts ([most] papers mark them; references/concepts.md) --
% ONE colour for every concept (house preference: a single orange).
% Colour is ON by default for drafting/review; \showconceptsfalse before the
% submission/camera-ready build renders everything black (house practice).
% Declare each concept once in the concept table of the main file
% ("% concept: <phrase> | core/framing | borrowed: <citekey> "<quote>" / composed: ... /
% coined | <definition>"; other key terms as "% term:" lines, concepts.md section 3), mark it
% with \concept{<phrase>} at first mention, and echo the SAME phrase later.
% Budget: 2-5 core (max 6) + <= 6 framing, 8-15 marked uses (cap 20),
% Introduction + first formal definition only. Borrow terms before coining.
\newif\ifshowconcepts
\showconceptstrue                 % \showconceptsfalse for submission
\definecolor{conceptcolor}{named}{orange}   % xcolor's orange, as in the drafting macros (B4)
\newcommand{\concept}[1]{\ifshowconcepts{\color{conceptcolor}#1}\else{#1}\fi}

\newcommand{\coB}[1]{{\color{black}{#1}}}     % senior co-author edits
\newcommand{\coC}[1]{{\color{black}{#1}}}
\newcommand{\possible}[1]{{\color{black}{#1}}}
\newcommand{\ignore}[1]{}                     % park deleted text inline [some]
\newcommand{\yes}[1]{{\color{teal}{#1}}}
\newcommand{\no}[1]{{\color{red}{#1}}}

% ---------- B5. Math kit ([all]; guarded against ICLR's math_commands.tex) ---
\ifdefined\argmax\else \DeclareMathOperator*{\argmax}{arg\,max}\fi
\ifdefined\argmin\else \DeclareMathOperator*{\argmin}{arg\,min}\fi
\ifdefined\sign\else   \DeclareMathOperator{\sign}{sign}\fi
\ifdefined\Tr\else     \DeclareMathOperator{\Tr}{Tr}\fi
\ifdefined\ab\else     \let\ab\allowbreak\fi
\providecommand\norm[1]{\left\lVert#1\right\rVert}
\providecommand{\Lapl}{\mathbf{\mathop{\mathcal{L}}}}   % losses: \Lapl_{\text{Name}}
\providecommand{\Trans}[1]{{#1}^{\top}}
\providecommand{\Bracs}[1]{\left({#1}\right)}
\providecommand{\Mat}[1]{\mathbf{#1}}
\providecommand{\Space}[1]{\mathbb{#1}}
\providecommand{\Set}[1]{\mathcal{#1}}
\providecommand{\ParaTitle}[1]{\noindent\textbf{#1}}
\providecommand\myeq{\mathrel{\stackrel{\makebox[0pt]{\mbox{\normalfont\tiny def}}}{=}}}
\providecommand\eqc{\mathrel{\overset{\makebox[0pt]{\mbox{\normalfont\tiny\sffamily c}}}{=}}}
% Definitions: use \triangleq [current].

% ---------- B6. Theorem environments (two numbering schemes, each [some]) ---
%  * Headline theorem: restatable and numbered within the section, printing
%    "Theorem 3.1"; \MainTheorem* restates it in the appendix. Use it when the
%    paper has ONE headline result (references/math.md section 2).
%  * Definitions and lemmas: plain counters with no section prefix, printing
%    "Definition 1" and "Lemma 1" [current]. Use them in every theory paper.
%  * The shared-counter set (theorem/proposition/lemma/... numbered within the
%    section) that the NeurIPS and ICML templates ship is NOT used here. ICML's
%    kit prefers it (a preference, not a rule): to follow it, replace this block
%    with the template's (references/venues.md section 3.4).
% Usage: \begin{restatable}{thm}{MainTheorem} ... \end{restatable}; appendix: \MainTheorem*
\declaretheorem[name=Theorem,numberwithin=section]{thm}        % headline, restatable
\newtheorem{definition}{Definition}                           % plain counter
\newtheorem{lemma}{Lemma}                                     % plain counter
% [skill choice]: the same plain counters for the remaining kinds, and a
% non-restatable "theorem" that shares thm's counter.
\newtheorem{theorem}[thm]{Theorem}
\newtheorem{proposition}{Proposition}
\newtheorem{corollary}{Corollary}
\newtheorem{assumption}{Assumption}
\newtheorem{remark}{Remark}
% For objectives that are derived rather than proved [some]:
\newenvironment{derivation_out}{\renewcommand{\proofname}{Derivation Outline}\proof}{\endproof}
\newenvironment{derivation_full}{\renewcommand{\proofname}{Full Derivation}\proof}{\endproof}

% ---------- B7. Colours ----------------------------------------------------
% (a) tables
\definecolor{+}{rgb}{0.70,0.13,0.13}          % #B32121 gain:  $^{\color{+}+1.15}$ [some]
\definecolor{-}{rgb}{0.25,0.41,0.88}          % #4069E0 drop:  $^{\color{-}-0.42}$ [some]
\definecolor{up}{rgb}{0.70,0.13,0.13}         % readable aliases
\definecolor{down}{rgb}{0.25,0.41,0.88}
\definecolor{oursrow}{gray}{0.95}             % #F2F2F2 our row (gray)
\definecolor{oursrowblue}{RGB}{239,246,251}   % #EFF6FB our row (blue; ICML-safe)
\definecolor{qualityrowcolor}{RGB}{255,248,249} % #FFF8F9 quality rows [some]
\definecolor{costrowcolor}{RGB}{246,249,254}     % #F6F9FE cost rows [some]
% (b) semantic text and case-study box colours [most]
\definecolor{good}{RGB}{58,113,104}
\definecolor{bad}{RGB}{180,0,0}
\definecolor{badboxborder}{RGB}{180,0,0}
\definecolor{goodboxborder}{RGB}{81,91,131}   % = accent-B outline in the diagrams
\definecolor{badboxbg}{RGB}{255,250,250}
\definecolor{goodboxbg}{RGB}{251,252,254}
\definecolor{titlebg}{RGB}{77,77,77}          % prompt-box title bar
% (c) role colours shared with diagrams and plots [current]; rename per paper
\definecolor{roleA}{RGB}{220,105,123}         % #DC697B (role A)
\definecolor{roleB}{RGB}{22,78,164}           % #164EA4 (role / mode B)
% (d) algorithm phase shading [current]
\definecolor{PhaseABg}{RGB}{234,208,228}      % e.g. Sampling
\definecolor{PhaseBBg}{RGB}{250,235,206}      % e.g. Scoring
\definecolor{PhaseCBg}{RGB}{217,215,219}      % e.g. Updating
% (e) plot palette as LaTeX colours (same hex as templates/plot_style.py)
\definecolor{plotblue}{HTML}{649CD8}
\definecolor{plotbluelight}{HTML}{BFD8EF}
\definecolor{plotpink}{HTML}{FF85A1}
\definecolor{plotpinklight}{HTML}{FFDAE2}
\definecolor{plotgray}{HTML}{595959}

% ---------- B8. Table helpers ----------------------------------------------
\newcommand{\pend}{--}                            % pending result cell (no data yet)
\newcommand{\ablation}[1]{{\color{gray}{#1}}}     % greyed ablation row [some]
% Set \renewcommand{\arraystretch}{0.9} INSIDE each table float, never globally.
% \shrinktofit{<max width>}{<tabular>} -- [skill choice].
% House practice writes \resizebox{w\linewidth}{!}{...}, which also ENLARGES a
% table narrower than w (its text then prints larger than body text). This
% wrapper scales only when the tabular is wider than <max width>, so wide
% tables look exactly like the house tables and narrow ones keep body-text size.
\newsavebox{\shrinkbox}
\newcommand{\shrinktofit}[2]{%
  \sbox{\shrinkbox}{#2}%
  \ifdim\wd\shrinkbox>#1\relax\resizebox{#1}{!}{\usebox{\shrinkbox}}%
  \else\usebox{\shrinkbox}\fi}

% ---------- B9. Figure placeholder (drafts without results) ----------------
% \figplaceholder[<height>]{<what the panel will show>}
\newcommand{\figplaceholder}[2][3cm]{%
  \fbox{\parbox[c][#1][c]{0.94\linewidth}{\centering\footnotesize\ttfamily #2}}}

% ---------- B10. Prompt and case-study boxes ([most]) ------------------------
\newtcolorbox{gptpromptbox}[1][]{colback=white,colframe=black,boxrule=1pt,arc=2pt,
  width=\linewidth,boxsep=2pt,left=2pt,right=2pt,top=2pt,bottom=2pt,
  title={PROMPT:},fonttitle=\bfseries,coltitle=white,colbacktitle=titlebg,#1}
\newtcolorbox{promptbox}[1][]{colback=white,colframe=black,boxrule=1pt,sharp corners,
  width=\linewidth,boxsep=2pt,left=2pt,right=2pt,top=2pt,bottom=2pt,#1}
\newtcolorbox{shadowpromptbox}[1][]{colback=white,colframe=black,boxrule=1pt,arc=2pt,
  width=\linewidth,boxsep=2pt,left=2pt,right=2pt,top=2pt,bottom=2pt,
  enhanced,drop fuzzy shadow=black!15!white,#1}                      % [current]
\newtcolorbox{badcasebox}{colback=badboxbg,colframe=badboxborder,sharp corners,
  boxrule=0.6pt,left=1pt,right=1pt,top=1pt,bottom=1pt}
\newtcolorbox{goodcasebox}{colback=goodboxbg,colframe=goodboxborder,sharp corners,
  boxrule=0.6pt,left=1pt,right=1pt,top=1pt,bottom=1pt}
\newsavebox{\PromptBoxA}
\newsavebox{\CaseStudy}

% ---------- B11. Algorithm helpers ([current]) ------------------------------
\newcommand{\CommentBelow}[1]{{\color{gray}\footnotesize // #1}}
\newcommand{\Cmt}[1]{\hfill{\color{gray}\footnotesize // #1}}
\providecommand{\BREAK}{\STATE \textbf{break}}
\pgfdeclarelayer{algobg}
\pgfsetlayers{algobg,main}
\newcommand{\tikzmk}[1]{\tikz[remember picture,overlay]\node[inner sep=0pt,outer sep=0pt](#1){};}
\newlength{\algoboxshiftx}\setlength{\algoboxshiftx}{-0.6em} % aligns in one and two columns (ICML skeleton);
                                                            % re-check if a mark sits at the end of a comment line
\newlength{\algoboxpadxL}\setlength{\algoboxpadxL}{0.6em}
\newlength{\algoboxpadxR}\setlength{\algoboxpadxR}{-1.2em}
\newlength{\algoboxpadyT}\setlength{\algoboxpadyT}{1.7ex}
\newlength{\algoboxpadyB}\setlength{\algoboxpadyB}{0.8ex}
% usage: \STATE \tikzmk{A}\CommentBelow{...} ... \STATE ...\tikzmk{B}
%        after \end{algorithmic}:  \boxit[Label]{PhaseABg}{A}{B}
% \boxit must sit OUTSIDE algorithmic; tikz marks need two compile passes.
\newcommand{\boxit}[4][]{%
\begin{tikzpicture}[remember picture,overlay]
\begin{pgfonlayer}{algobg}
  \coordinate (ALG@NW) at ([xshift=\algoboxshiftx-\algoboxpadxL,yshift=\algoboxpadyT]#3);
  \coordinate (ALG@SE) at ([xshift=\algoboxshiftx+\dimexpr\linewidth+\algoboxpadxR\relax,
                            yshift=-\algoboxpadyB]#3 |- #4);
  \fill[#2, opacity=0.14, rounded corners=2pt] (ALG@NW) rectangle (ALG@SE);
  \if\relax\detokenize{#1}\relax\else
    \node[anchor=south east, font=\scriptsize\bfseries, text=black!70]
      at ([xshift=-0.2em,yshift=0.2ex]ALG@SE) {#1};
  \fi
\end{pgfonlayer}
\end{tikzpicture}%
}
 after math_commands.tex is equivalent.
%%
%% (A2) NeurIPS 2026
%%   \documentclass{article}
%%   \PassOptionsToPackage{table}{xcolor}
%%   \PassOptionsToPackage{numbers, compress}{natbib}
%%   \usepackage[main]{neurips_2026}        % camera-ready: [main,final]
%%                                          % ([final] alone stops at \@trackname)
%%   \newcommand{\VenueLinks}{\hypersetup{colorlinks, linkcolor=red,
%%       anchorcolor=blue, citecolor=green}}
%%   Cite with \cite{} (numeric). \bibliographystyle{unsrtnat}.
%%   After the appendix: \clearpage \input{checklist}. 9 content pages.
%%
%% (A3) ICML 2026  [full main file: main_skeleton_icml.tex]
%%   \documentclass{article}
%%   \PassOptionsToPackage{table}{xcolor}   % prevents "Option clash for package
%%                                          % xcolor" when [table]{xcolor} follows the sty
%%   \usepackage[T1]{fontenc}               % B1 loads it too; without T1, XeTeX
%%                                          % silently falls back to Latin Modern
%%   \usepackage{microtype}\usepackage{graphicx}\usepackage{subcaption}
%%   \usepackage{booktabs}\usepackage{hyperref}
%%   \newcommand{\theHalgorithm}{\arabic{algorithm}}
%%   \usepackage{icml2026}                  % blind submission (line-number ruler);
%%                                          % [preprint] arXiv; [accepted] camera-ready
%%   \icmltitlerunning{Short Running Title}
%%   After the \twocolumn[...] title block ALWAYS call \printAffiliationsAndNotice{}
%%   (official template); without it icml2026.sty's end-of-document warning fails
%%   with "Unbalanced write command" (verified). Only [accepted] needs CM italic 9pt
%%   (cmti9); if tectonic times out fetching it, use scripts/tectonic_fetch.sh and
%%   -Z search-path=<dir> (references/venues.md section 2.3).
%%   No \VenueLinks (keep the template's dark-blue links). \cite{} is fine: the sty
%%   maps it to \citep{} (parenthetical, ONE argument; \citep[p.~3]{x} for pages).
%%   8 pp at submission (+1 camera-ready); \section*{Impact Statement} (required)
%%   after \clearpage, before the references. Plots use \columnwidth. The kit is
%%   not bundled: templates/venue/icml/README.md.
%% =====================================================================

% ---------- B1. Encoding --------------------------------------------------
% T1 under BOTH engines: the venue styles select Times ("ptm"), which has no
% TU shape; without T1, XeLaTeX falls back to Latin Modern.
\usepackage{iftex}
\ifPDFTeX
  \usepackage[utf8]{inputenc}
\fi
\usepackage[T1]{fontenc}

% ---------- B2. Packages ([all] core + [current] additions) ---------------
\usepackage{xcolor}                 % [table] comes from \PassOptionsToPackage in Part A
\usepackage{url}
\usepackage{booktabs}
\usepackage{amsmath,amssymb,amsfonts}
\usepackage{mathtools}
\usepackage{amsthm}
\usepackage{thmtools}               % provides thm-restate (restatable theorems)
% no bm [current]: write \mathbf{} (\boldsymbol{} from amsmath if needed); with
% bm the NeurIPS 8pt footnotes need cmmib8, which the offline tectonic cache
% lacked (build stopped; verified).
\usepackage{nicefrac}
\usepackage{microtype}
\usepackage{graphicx}
\usepackage{caption}
\usepackage{subcaption}
\usepackage{enumitem}               % always \begin{itemize}[leftmargin=*]
\usepackage{multirow}
\usepackage{tabularx}
\usepackage{makecell}
\usepackage{wrapfig}                % NeurIPS/ICLR only; never in two-column ICML
\usepackage{tcolorbox}
\tcbuselibrary{skins}
\usepackage{tikz}
\usetikzlibrary{calc,shadows}
\ifdefined\algorithmic\else         % ICML's icml2026.sty already loads these
  \usepackage{algorithm}
  \usepackage{algorithmic}
\fi
\ifdefined\hypersetup\else \usepackage{hyperref}\fi
\ifdefined\VenueLinks \VenueLinks \fi
% No cleveref: write "Figure~\ref{}", "Table~\ref{}", "Section~\ref{}",
% "Appendix~\ref{}", "Equation~\ref{}" by hand (house practice types the word
% by hand [all]; the tie only blocks a line break before the number:
% references/latex_source.md section 3). No method-name macro: type the
% name [all].

% ---------- B3. Latin abbreviations ([all]; trailing space is intentional) ---
% Use inside parentheses: (\ie ...), (\eg ...), (\cf Section~\ref{...}).
% Do not type another comma or space after them.
\providecommand{\ie}{\emph{i.e., }}
\providecommand{\eg}{\emph{e.g., }}
\providecommand{\etal}{\emph{et al.}}
\providecommand{\st}{\emph{s.t. }}
\providecommand{\etc}{\emph{etc.}}
\providecommand{\wrt}{\emph{w.r.t.}~}
\providecommand{\cf}{\emph{cf. }}
\providecommand{\aka}{\emph{aka. }}

% ---------- B4. Co-author / status macros ([all]) ---------------------------
% One macro per co-author, named after that person's initials: rename \coA,
% \coB, \coC to each co-author's initials. Coloured while drafting, black for
% submission and camera-ready; the wrappers stay in the text. House practice
% swaps two commented blocks; here one switch does the same.
\newif\ifdraftcolors
\draftcolorsfalse                  % \draftcolorstrue while drafting
\ifdraftcolors
  \newcommand{\coA}[1]{{\color{orange}{#1}}}
  \newcommand{\revision}[1]{{\color{blue}{#1}}}
  \newcommand{\rebuttal}[1]{{\color{blue}{#1}}}
  \newcommand{\draft}[1]{{\color{purple}{#1}}}
  \newcommand{\warning}[1]{{\color{red}{#1}}}
\else
  \newcommand{\coA}[1]{{\color{black}{#1}}}
  \newcommand{\revision}[1]{{\color{black}{#1}}}
  \newcommand{\rebuttal}[1]{{\color{black}{#1}}}
  \newcommand{\draft}[1]{{\color{black}{#1}}}
  \newcommand{\warning}[1]{{\color{black}{#1}}}
\fi
% ---------- B4b. Key concepts ([most] papers mark them; references/concepts.md) --
% ONE colour for every concept (house preference: a single orange).
% Colour is ON by default for drafting/review; \showconceptsfalse before the
% submission/camera-ready build renders everything black (house practice).
% Declare each concept once in the concept table of the main file
% ("% concept: <phrase> | core/framing | borrowed: <citekey> "<quote>" / composed: ... /
% coined | <definition>"; other key terms as "% term:" lines, concepts.md section 3), mark it
% with \concept{<phrase>} at first mention, and echo the SAME phrase later.
% Budget: 2-5 core (max 6) + <= 6 framing, 8-15 marked uses (cap 20),
% Introduction + first formal definition only. Borrow terms before coining.
\newif\ifshowconcepts
\showconceptstrue                 % \showconceptsfalse for submission
\definecolor{conceptcolor}{named}{orange}   % xcolor's orange, as in the drafting macros (B4)
\newcommand{\concept}[1]{\ifshowconcepts{\color{conceptcolor}#1}\else{#1}\fi}

\newcommand{\coB}[1]{{\color{black}{#1}}}     % senior co-author edits
\newcommand{\coC}[1]{{\color{black}{#1}}}
\newcommand{\possible}[1]{{\color{black}{#1}}}
\newcommand{\ignore}[1]{}                     % park deleted text inline [some]
\newcommand{\yes}[1]{{\color{teal}{#1}}}
\newcommand{\no}[1]{{\color{red}{#1}}}

% ---------- B5. Math kit ([all]; guarded against ICLR's math_commands.tex) ---
\ifdefined\argmax\else \DeclareMathOperator*{\argmax}{arg\,max}\fi
\ifdefined\argmin\else \DeclareMathOperator*{\argmin}{arg\,min}\fi
\ifdefined\sign\else   \DeclareMathOperator{\sign}{sign}\fi
\ifdefined\Tr\else     \DeclareMathOperator{\Tr}{Tr}\fi
\ifdefined\ab\else     \let\ab\allowbreak\fi
\providecommand\norm[1]{\left\lVert#1\right\rVert}
\providecommand{\Lapl}{\mathbf{\mathop{\mathcal{L}}}}   % losses: \Lapl_{\text{Name}}
\providecommand{\Trans}[1]{{#1}^{\top}}
\providecommand{\Bracs}[1]{\left({#1}\right)}
\providecommand{\Mat}[1]{\mathbf{#1}}
\providecommand{\Space}[1]{\mathbb{#1}}
\providecommand{\Set}[1]{\mathcal{#1}}
\providecommand{\ParaTitle}[1]{\noindent\textbf{#1}}
\providecommand\myeq{\mathrel{\stackrel{\makebox[0pt]{\mbox{\normalfont\tiny def}}}{=}}}
\providecommand\eqc{\mathrel{\overset{\makebox[0pt]{\mbox{\normalfont\tiny\sffamily c}}}{=}}}
% Definitions: use \triangleq [current].

% ---------- B6. Theorem environments (two numbering schemes, each [some]) ---
%  * Headline theorem: restatable and numbered within the section, printing
%    "Theorem 3.1"; \MainTheorem* restates it in the appendix. Use it when the
%    paper has ONE headline result (references/math.md section 2).
%  * Definitions and lemmas: plain counters with no section prefix, printing
%    "Definition 1" and "Lemma 1" [current]. Use them in every theory paper.
%  * The shared-counter set (theorem/proposition/lemma/... numbered within the
%    section) that the NeurIPS and ICML templates ship is NOT used here. ICML's
%    kit prefers it (a preference, not a rule): to follow it, replace this block
%    with the template's (references/venues.md section 3.4).
% Usage: \begin{restatable}{thm}{MainTheorem} ... \end{restatable}; appendix: \MainTheorem*
\declaretheorem[name=Theorem,numberwithin=section]{thm}        % headline, restatable
\newtheorem{definition}{Definition}                           % plain counter
\newtheorem{lemma}{Lemma}                                     % plain counter
% [skill choice]: the same plain counters for the remaining kinds, and a
% non-restatable "theorem" that shares thm's counter.
\newtheorem{theorem}[thm]{Theorem}
\newtheorem{proposition}{Proposition}
\newtheorem{corollary}{Corollary}
\newtheorem{assumption}{Assumption}
\newtheorem{remark}{Remark}
% For objectives that are derived rather than proved [some]:
\newenvironment{derivation_out}{\renewcommand{\proofname}{Derivation Outline}\proof}{\endproof}
\newenvironment{derivation_full}{\renewcommand{\proofname}{Full Derivation}\proof}{\endproof}

% ---------- B7. Colours ----------------------------------------------------
% (a) tables
\definecolor{+}{rgb}{0.70,0.13,0.13}          % #B32121 gain:  $^{\color{+}+1.15}$ [some]
\definecolor{-}{rgb}{0.25,0.41,0.88}          % #4069E0 drop:  $^{\color{-}-0.42}$ [some]
\definecolor{up}{rgb}{0.70,0.13,0.13}         % readable aliases
\definecolor{down}{rgb}{0.25,0.41,0.88}
\definecolor{oursrow}{gray}{0.95}             % #F2F2F2 our row (gray)
\definecolor{oursrowblue}{RGB}{239,246,251}   % #EFF6FB our row (blue; ICML-safe)
\definecolor{qualityrowcolor}{RGB}{255,248,249} % #FFF8F9 quality rows [some]
\definecolor{costrowcolor}{RGB}{246,249,254}     % #F6F9FE cost rows [some]
% (b) semantic text and case-study box colours [most]
\definecolor{good}{RGB}{58,113,104}
\definecolor{bad}{RGB}{180,0,0}
\definecolor{badboxborder}{RGB}{180,0,0}
\definecolor{goodboxborder}{RGB}{81,91,131}   % = accent-B outline in the diagrams
\definecolor{badboxbg}{RGB}{255,250,250}
\definecolor{goodboxbg}{RGB}{251,252,254}
\definecolor{titlebg}{RGB}{77,77,77}          % prompt-box title bar
% (c) role colours shared with diagrams and plots [current]; rename per paper
\definecolor{roleA}{RGB}{220,105,123}         % #DC697B (role A)
\definecolor{roleB}{RGB}{22,78,164}           % #164EA4 (role / mode B)
% (d) algorithm phase shading [current]
\definecolor{PhaseABg}{RGB}{234,208,228}      % e.g. Sampling
\definecolor{PhaseBBg}{RGB}{250,235,206}      % e.g. Scoring
\definecolor{PhaseCBg}{RGB}{217,215,219}      % e.g. Updating
% (e) plot palette as LaTeX colours (same hex as templates/plot_style.py)
\definecolor{plotblue}{HTML}{649CD8}
\definecolor{plotbluelight}{HTML}{BFD8EF}
\definecolor{plotpink}{HTML}{FF85A1}
\definecolor{plotpinklight}{HTML}{FFDAE2}
\definecolor{plotgray}{HTML}{595959}

% ---------- B8. Table helpers ----------------------------------------------
\newcommand{\pend}{--}                            % pending result cell (no data yet)
\newcommand{\ablation}[1]{{\color{gray}{#1}}}     % greyed ablation row [some]
% Set \renewcommand{\arraystretch}{0.9} INSIDE each table float, never globally.
% \shrinktofit{<max width>}{<tabular>} -- [skill choice].
% House practice writes \resizebox{w\linewidth}{!}{...}, which also ENLARGES a
% table narrower than w (its text then prints larger than body text). This
% wrapper scales only when the tabular is wider than <max width>, so wide
% tables look exactly like the house tables and narrow ones keep body-text size.
\newsavebox{\shrinkbox}
\newcommand{\shrinktofit}[2]{%
  \sbox{\shrinkbox}{#2}%
  \ifdim\wd\shrinkbox>#1\relax\resizebox{#1}{!}{\usebox{\shrinkbox}}%
  \else\usebox{\shrinkbox}\fi}

% ---------- B9. Figure placeholder (drafts without results) ----------------
% \figplaceholder[<height>]{<what the panel will show>}
\newcommand{\figplaceholder}[2][3cm]{%
  \fbox{\parbox[c][#1][c]{0.94\linewidth}{\centering\footnotesize\ttfamily #2}}}

% ---------- B10. Prompt and case-study boxes ([most]) ------------------------
\newtcolorbox{gptpromptbox}[1][]{colback=white,colframe=black,boxrule=1pt,arc=2pt,
  width=\linewidth,boxsep=2pt,left=2pt,right=2pt,top=2pt,bottom=2pt,
  title={PROMPT:},fonttitle=\bfseries,coltitle=white,colbacktitle=titlebg,#1}
\newtcolorbox{promptbox}[1][]{colback=white,colframe=black,boxrule=1pt,sharp corners,
  width=\linewidth,boxsep=2pt,left=2pt,right=2pt,top=2pt,bottom=2pt,#1}
\newtcolorbox{shadowpromptbox}[1][]{colback=white,colframe=black,boxrule=1pt,arc=2pt,
  width=\linewidth,boxsep=2pt,left=2pt,right=2pt,top=2pt,bottom=2pt,
  enhanced,drop fuzzy shadow=black!15!white,#1}                      % [current]
\newtcolorbox{badcasebox}{colback=badboxbg,colframe=badboxborder,sharp corners,
  boxrule=0.6pt,left=1pt,right=1pt,top=1pt,bottom=1pt}
\newtcolorbox{goodcasebox}{colback=goodboxbg,colframe=goodboxborder,sharp corners,
  boxrule=0.6pt,left=1pt,right=1pt,top=1pt,bottom=1pt}
\newsavebox{\PromptBoxA}
\newsavebox{\CaseStudy}

% ---------- B11. Algorithm helpers ([current]) ------------------------------
\newcommand{\CommentBelow}[1]{{\color{gray}\footnotesize // #1}}
\newcommand{\Cmt}[1]{\hfill{\color{gray}\footnotesize // #1}}
\providecommand{\BREAK}{\STATE \textbf{break}}
\pgfdeclarelayer{algobg}
\pgfsetlayers{algobg,main}
\newcommand{\tikzmk}[1]{\tikz[remember picture,overlay]\node[inner sep=0pt,outer sep=0pt](#1){};}
\newlength{\algoboxshiftx}\setlength{\algoboxshiftx}{-0.6em} % aligns in one and two columns (ICML skeleton);
                                                            % re-check if a mark sits at the end of a comment line
\newlength{\algoboxpadxL}\setlength{\algoboxpadxL}{0.6em}
\newlength{\algoboxpadxR}\setlength{\algoboxpadxR}{-1.2em}
\newlength{\algoboxpadyT}\setlength{\algoboxpadyT}{1.7ex}
\newlength{\algoboxpadyB}\setlength{\algoboxpadyB}{0.8ex}
% usage: \STATE \tikzmk{A}\CommentBelow{...} ... \STATE ...\tikzmk{B}
%        after \end{algorithmic}:  \boxit[Label]{PhaseABg}{A}{B}
% \boxit must sit OUTSIDE algorithmic; tikz marks need two compile passes.
\newcommand{\boxit}[4][]{%
\begin{tikzpicture}[remember picture,overlay]
\begin{pgfonlayer}{algobg}
  \coordinate (ALG@NW) at ([xshift=\algoboxshiftx-\algoboxpadxL,yshift=\algoboxpadyT]#3);
  \coordinate (ALG@SE) at ([xshift=\algoboxshiftx+\dimexpr\linewidth+\algoboxpadxR\relax,
                            yshift=-\algoboxpadyB]#3 |- #4);
  \fill[#2, opacity=0.14, rounded corners=2pt] (ALG@NW) rectangle (ALG@SE);
  \if\relax\detokenize{#1}\relax\else
    \node[anchor=south east, font=\scriptsize\bfseries, text=black!70]
      at ([xshift=-0.2em,yshift=0.2ex]ALG@SE) {#1};
  \fi
\end{pgfonlayer}
\end{tikzpicture}%
}
 after math_commands.tex is equivalent.
%%
%% (A2) NeurIPS 2026
%%   \documentclass{article}
%%   \PassOptionsToPackage{table}{xcolor}
%%   \PassOptionsToPackage{numbers, compress}{natbib}
%%   \usepackage[main]{neurips_2026}        % camera-ready: [main,final]
%%                                          % ([final] alone stops at \@trackname)
%%   \newcommand{\VenueLinks}{\hypersetup{colorlinks, linkcolor=red,
%%       anchorcolor=blue, citecolor=green}}
%%   Cite with \cite{} (numeric). \bibliographystyle{unsrtnat}.
%%   After the appendix: \clearpage \input{checklist}. 9 content pages.
%%
%% (A3) ICML 2026  [full main file: main_skeleton_icml.tex]
%%   \documentclass{article}
%%   \PassOptionsToPackage{table}{xcolor}   % prevents "Option clash for package
%%                                          % xcolor" when [table]{xcolor} follows the sty
%%   \usepackage[T1]{fontenc}               % B1 loads it too; without T1, XeTeX
%%                                          % silently falls back to Latin Modern
%%   \usepackage{microtype}\usepackage{graphicx}\usepackage{subcaption}
%%   \usepackage{booktabs}\usepackage{hyperref}
%%   \newcommand{\theHalgorithm}{\arabic{algorithm}}
%%   \usepackage{icml2026}                  % blind submission (line-number ruler);
%%                                          % [preprint] arXiv; [accepted] camera-ready
%%   \icmltitlerunning{Short Running Title}
%%   After the \twocolumn[...] title block ALWAYS call \printAffiliationsAndNotice{}
%%   (official template); without it icml2026.sty's end-of-document warning fails
%%   with "Unbalanced write command" (verified). Only [accepted] needs CM italic 9pt
%%   (cmti9); if tectonic times out fetching it, use scripts/tectonic_fetch.sh and
%%   -Z search-path=<dir> (references/venues.md section 2.3).
%%   No \VenueLinks (keep the template's dark-blue links). \cite{} is fine: the sty
%%   maps it to \citep{} (parenthetical, ONE argument; \citep[p.~3]{x} for pages).
%%   8 pp at submission (+1 camera-ready); \section*{Impact Statement} (required)
%%   after \clearpage, before the references. Plots use \columnwidth. The kit is
%%   not bundled: templates/venue/icml/README.md.
%% =====================================================================

% ---------- B1. Encoding --------------------------------------------------
% T1 under BOTH engines: the venue styles select Times ("ptm"), which has no
% TU shape; without T1, XeLaTeX falls back to Latin Modern.
\usepackage{iftex}
\ifPDFTeX
  \usepackage[utf8]{inputenc}
\fi
\usepackage[T1]{fontenc}

% ---------- B2. Packages ([all] core + [current] additions) ---------------
\usepackage{xcolor}                 % [table] comes from \PassOptionsToPackage in Part A
\usepackage{url}
\usepackage{booktabs}
\usepackage{amsmath,amssymb,amsfonts}
\usepackage{mathtools}
\usepackage{amsthm}
\usepackage{thmtools}               % provides thm-restate (restatable theorems)
% no bm [current]: write \mathbf{} (\boldsymbol{} from amsmath if needed); with
% bm the NeurIPS 8pt footnotes need cmmib8, which the offline tectonic cache
% lacked (build stopped; verified).
\usepackage{nicefrac}
\usepackage{microtype}
\usepackage{graphicx}
\usepackage{caption}
\usepackage{subcaption}
\usepackage{enumitem}               % always \begin{itemize}[leftmargin=*]
\usepackage{multirow}
\usepackage{tabularx}
\usepackage{makecell}
\usepackage{wrapfig}                % NeurIPS/ICLR only; never in two-column ICML
\usepackage{tcolorbox}
\tcbuselibrary{skins}
\usepackage{tikz}
\usetikzlibrary{calc,shadows}
\ifdefined\algorithmic\else         % ICML's icml2026.sty already loads these
  \usepackage{algorithm}
  \usepackage{algorithmic}
\fi
\ifdefined\hypersetup\else \usepackage{hyperref}\fi
\ifdefined\VenueLinks \VenueLinks \fi
% No cleveref: write "Figure~\ref{}", "Table~\ref{}", "Section~\ref{}",
% "Appendix~\ref{}", "Equation~\ref{}" by hand (house practice types the word
% by hand [all]; the tie only blocks a line break before the number:
% references/latex_source.md section 3). No method-name macro: type the
% name [all].

% ---------- B3. Latin abbreviations ([all]; trailing space is intentional) ---
% Use inside parentheses: (\ie ...), (\eg ...), (\cf Section~\ref{...}).
% Do not type another comma or space after them.
\providecommand{\ie}{\emph{i.e., }}
\providecommand{\eg}{\emph{e.g., }}
\providecommand{\etal}{\emph{et al.}}
\providecommand{\st}{\emph{s.t. }}
\providecommand{\etc}{\emph{etc.}}
\providecommand{\wrt}{\emph{w.r.t.}~}
\providecommand{\cf}{\emph{cf. }}
\providecommand{\aka}{\emph{aka. }}

% ---------- B4. Co-author / status macros ([all]) ---------------------------
% One macro per co-author, named after that person's initials: rename \coA,
% \coB, \coC to each co-author's initials. Coloured while drafting, black for
% submission and camera-ready; the wrappers stay in the text. House practice
% swaps two commented blocks; here one switch does the same.
\newif\ifdraftcolors
\draftcolorsfalse                  % \draftcolorstrue while drafting
\ifdraftcolors
  \newcommand{\coA}[1]{{\color{orange}{#1}}}
  \newcommand{\revision}[1]{{\color{blue}{#1}}}
  \newcommand{\rebuttal}[1]{{\color{blue}{#1}}}
  \newcommand{\draft}[1]{{\color{purple}{#1}}}
  \newcommand{\warning}[1]{{\color{red}{#1}}}
\else
  \newcommand{\coA}[1]{{\color{black}{#1}}}
  \newcommand{\revision}[1]{{\color{black}{#1}}}
  \newcommand{\rebuttal}[1]{{\color{black}{#1}}}
  \newcommand{\draft}[1]{{\color{black}{#1}}}
  \newcommand{\warning}[1]{{\color{black}{#1}}}
\fi
% ---------- B4b. Key concepts ([most] papers mark them; references/concepts.md) --
% ONE colour for every concept (house preference: a single orange).
% Colour is ON by default for drafting/review; \showconceptsfalse before the
% submission/camera-ready build renders everything black (house practice).
% Declare each concept once in the concept table of the main file
% ("% concept: <phrase> | core/framing | borrowed: <citekey> "<quote>" / composed: ... /
% coined | <definition>"; other key terms as "% term:" lines, concepts.md section 3), mark it
% with \concept{<phrase>} at first mention, and echo the SAME phrase later.
% Budget: 2-5 core (max 6) + <= 6 framing, 8-15 marked uses (cap 20),
% Introduction + first formal definition only. Borrow terms before coining.
\newif\ifshowconcepts
\showconceptstrue                 % \showconceptsfalse for submission
\definecolor{conceptcolor}{named}{orange}   % xcolor's orange, as in the drafting macros (B4)
\newcommand{\concept}[1]{\ifshowconcepts{\color{conceptcolor}#1}\else{#1}\fi}

\newcommand{\coB}[1]{{\color{black}{#1}}}     % senior co-author edits
\newcommand{\coC}[1]{{\color{black}{#1}}}
\newcommand{\possible}[1]{{\color{black}{#1}}}
\newcommand{\ignore}[1]{}                     % park deleted text inline [some]
\newcommand{\yes}[1]{{\color{teal}{#1}}}
\newcommand{\no}[1]{{\color{red}{#1}}}

% ---------- B5. Math kit ([all]; guarded against ICLR's math_commands.tex) ---
\ifdefined\argmax\else \DeclareMathOperator*{\argmax}{arg\,max}\fi
\ifdefined\argmin\else \DeclareMathOperator*{\argmin}{arg\,min}\fi
\ifdefined\sign\else   \DeclareMathOperator{\sign}{sign}\fi
\ifdefined\Tr\else     \DeclareMathOperator{\Tr}{Tr}\fi
\ifdefined\ab\else     \let\ab\allowbreak\fi
\providecommand\norm[1]{\left\lVert#1\right\rVert}
\providecommand{\Lapl}{\mathbf{\mathop{\mathcal{L}}}}   % losses: \Lapl_{\text{Name}}
\providecommand{\Trans}[1]{{#1}^{\top}}
\providecommand{\Bracs}[1]{\left({#1}\right)}
\providecommand{\Mat}[1]{\mathbf{#1}}
\providecommand{\Space}[1]{\mathbb{#1}}
\providecommand{\Set}[1]{\mathcal{#1}}
\providecommand{\ParaTitle}[1]{\noindent\textbf{#1}}
\providecommand\myeq{\mathrel{\stackrel{\makebox[0pt]{\mbox{\normalfont\tiny def}}}{=}}}
\providecommand\eqc{\mathrel{\overset{\makebox[0pt]{\mbox{\normalfont\tiny\sffamily c}}}{=}}}
% Definitions: use \triangleq [current].

% ---------- B6. Theorem environments (two numbering schemes, each [some]) ---
%  * Headline theorem: restatable and numbered within the section, printing
%    "Theorem 3.1"; \MainTheorem* restates it in the appendix. Use it when the
%    paper has ONE headline result (references/math.md section 2).
%  * Definitions and lemmas: plain counters with no section prefix, printing
%    "Definition 1" and "Lemma 1" [current]. Use them in every theory paper.
%  * The shared-counter set (theorem/proposition/lemma/... numbered within the
%    section) that the NeurIPS and ICML templates ship is NOT used here. ICML's
%    kit prefers it (a preference, not a rule): to follow it, replace this block
%    with the template's (references/venues.md section 3.4).
% Usage: \begin{restatable}{thm}{MainTheorem} ... \end{restatable}; appendix: \MainTheorem*
\declaretheorem[name=Theorem,numberwithin=section]{thm}        % headline, restatable
\newtheorem{definition}{Definition}                           % plain counter
\newtheorem{lemma}{Lemma}                                     % plain counter
% [skill choice]: the same plain counters for the remaining kinds, and a
% non-restatable "theorem" that shares thm's counter.
\newtheorem{theorem}[thm]{Theorem}
\newtheorem{proposition}{Proposition}
\newtheorem{corollary}{Corollary}
\newtheorem{assumption}{Assumption}
\newtheorem{remark}{Remark}
% For objectives that are derived rather than proved [some]:
\newenvironment{derivation_out}{\renewcommand{\proofname}{Derivation Outline}\proof}{\endproof}
\newenvironment{derivation_full}{\renewcommand{\proofname}{Full Derivation}\proof}{\endproof}

% ---------- B7. Colours ----------------------------------------------------
% (a) tables
\definecolor{+}{rgb}{0.70,0.13,0.13}          % #B32121 gain:  $^{\color{+}+1.15}$ [some]
\definecolor{-}{rgb}{0.25,0.41,0.88}          % #4069E0 drop:  $^{\color{-}-0.42}$ [some]
\definecolor{up}{rgb}{0.70,0.13,0.13}         % readable aliases
\definecolor{down}{rgb}{0.25,0.41,0.88}
\definecolor{oursrow}{gray}{0.95}             % #F2F2F2 our row (gray)
\definecolor{oursrowblue}{RGB}{239,246,251}   % #EFF6FB our row (blue; ICML-safe)
\definecolor{qualityrowcolor}{RGB}{255,248,249} % #FFF8F9 quality rows [some]
\definecolor{costrowcolor}{RGB}{246,249,254}     % #F6F9FE cost rows [some]
% (b) semantic text and case-study box colours [most]
\definecolor{good}{RGB}{58,113,104}
\definecolor{bad}{RGB}{180,0,0}
\definecolor{badboxborder}{RGB}{180,0,0}
\definecolor{goodboxborder}{RGB}{81,91,131}   % = accent-B outline in the diagrams
\definecolor{badboxbg}{RGB}{255,250,250}
\definecolor{goodboxbg}{RGB}{251,252,254}
\definecolor{titlebg}{RGB}{77,77,77}          % prompt-box title bar
% (c) role colours shared with diagrams and plots [current]; rename per paper
\definecolor{roleA}{RGB}{220,105,123}         % #DC697B (role A)
\definecolor{roleB}{RGB}{22,78,164}           % #164EA4 (role / mode B)
% (d) algorithm phase shading [current]
\definecolor{PhaseABg}{RGB}{234,208,228}      % e.g. Sampling
\definecolor{PhaseBBg}{RGB}{250,235,206}      % e.g. Scoring
\definecolor{PhaseCBg}{RGB}{217,215,219}      % e.g. Updating
% (e) plot palette as LaTeX colours (same hex as templates/plot_style.py)
\definecolor{plotblue}{HTML}{649CD8}
\definecolor{plotbluelight}{HTML}{BFD8EF}
\definecolor{plotpink}{HTML}{FF85A1}
\definecolor{plotpinklight}{HTML}{FFDAE2}
\definecolor{plotgray}{HTML}{595959}

% ---------- B8. Table helpers ----------------------------------------------
\newcommand{\pend}{--}                            % pending result cell (no data yet)
\newcommand{\ablation}[1]{{\color{gray}{#1}}}     % greyed ablation row [some]
% Set \renewcommand{\arraystretch}{0.9} INSIDE each table float, never globally.
% \shrinktofit{<max width>}{<tabular>} -- [skill choice].
% House practice writes \resizebox{w\linewidth}{!}{...}, which also ENLARGES a
% table narrower than w (its text then prints larger than body text). This
% wrapper scales only when the tabular is wider than <max width>, so wide
% tables look exactly like the house tables and narrow ones keep body-text size.
\newsavebox{\shrinkbox}
\newcommand{\shrinktofit}[2]{%
  \sbox{\shrinkbox}{#2}%
  \ifdim\wd\shrinkbox>#1\relax\resizebox{#1}{!}{\usebox{\shrinkbox}}%
  \else\usebox{\shrinkbox}\fi}

% ---------- B9. Figure placeholder (drafts without results) ----------------
% \figplaceholder[<height>]{<what the panel will show>}
\newcommand{\figplaceholder}[2][3cm]{%
  \fbox{\parbox[c][#1][c]{0.94\linewidth}{\centering\footnotesize\ttfamily #2}}}

% ---------- B10. Prompt and case-study boxes ([most]) ------------------------
\newtcolorbox{gptpromptbox}[1][]{colback=white,colframe=black,boxrule=1pt,arc=2pt,
  width=\linewidth,boxsep=2pt,left=2pt,right=2pt,top=2pt,bottom=2pt,
  title={PROMPT:},fonttitle=\bfseries,coltitle=white,colbacktitle=titlebg,#1}
\newtcolorbox{promptbox}[1][]{colback=white,colframe=black,boxrule=1pt,sharp corners,
  width=\linewidth,boxsep=2pt,left=2pt,right=2pt,top=2pt,bottom=2pt,#1}
\newtcolorbox{shadowpromptbox}[1][]{colback=white,colframe=black,boxrule=1pt,arc=2pt,
  width=\linewidth,boxsep=2pt,left=2pt,right=2pt,top=2pt,bottom=2pt,
  enhanced,drop fuzzy shadow=black!15!white,#1}                      % [current]
\newtcolorbox{badcasebox}{colback=badboxbg,colframe=badboxborder,sharp corners,
  boxrule=0.6pt,left=1pt,right=1pt,top=1pt,bottom=1pt}
\newtcolorbox{goodcasebox}{colback=goodboxbg,colframe=goodboxborder,sharp corners,
  boxrule=0.6pt,left=1pt,right=1pt,top=1pt,bottom=1pt}
\newsavebox{\PromptBoxA}
\newsavebox{\CaseStudy}

% ---------- B11. Algorithm helpers ([current]) ------------------------------
\newcommand{\CommentBelow}[1]{{\color{gray}\footnotesize // #1}}
\newcommand{\Cmt}[1]{\hfill{\color{gray}\footnotesize // #1}}
\providecommand{\BREAK}{\STATE \textbf{break}}
\pgfdeclarelayer{algobg}
\pgfsetlayers{algobg,main}
\newcommand{\tikzmk}[1]{\tikz[remember picture,overlay]\node[inner sep=0pt,outer sep=0pt](#1){};}
\newlength{\algoboxshiftx}\setlength{\algoboxshiftx}{-0.6em} % aligns in one and two columns (ICML skeleton);
                                                            % re-check if a mark sits at the end of a comment line
\newlength{\algoboxpadxL}\setlength{\algoboxpadxL}{0.6em}
\newlength{\algoboxpadxR}\setlength{\algoboxpadxR}{-1.2em}
\newlength{\algoboxpadyT}\setlength{\algoboxpadyT}{1.7ex}
\newlength{\algoboxpadyB}\setlength{\algoboxpadyB}{0.8ex}
% usage: \STATE \tikzmk{A}\CommentBelow{...} ... \STATE ...\tikzmk{B}
%        after \end{algorithmic}:  \boxit[Label]{PhaseABg}{A}{B}
% \boxit must sit OUTSIDE algorithmic; tikz marks need two compile passes.
\newcommand{\boxit}[4][]{%
\begin{tikzpicture}[remember picture,overlay]
\begin{pgfonlayer}{algobg}
  \coordinate (ALG@NW) at ([xshift=\algoboxshiftx-\algoboxpadxL,yshift=\algoboxpadyT]#3);
  \coordinate (ALG@SE) at ([xshift=\algoboxshiftx+\dimexpr\linewidth+\algoboxpadxR\relax,
                            yshift=-\algoboxpadyB]#3 |- #4);
  \fill[#2, opacity=0.14, rounded corners=2pt] (ALG@NW) rectangle (ALG@SE);
  \if\relax\detokenize{#1}\relax\else
    \node[anchor=south east, font=\scriptsize\bfseries, text=black!70]
      at ([xshift=-0.2em,yshift=0.2ex]ALG@SE) {#1};
  \fi
\end{pgfonlayer}
\end{tikzpicture}%
}


% ---- Concept table (references/concepts.md): 2-5 core (max 6) + <= 6 framing terms.
% Borrow the literature's term before coining; one concept = one exact phrase; echo core concepts
% (same phrase, unmarked) in the Methodology, Experiments and Conclusion.
% Term table (references/concepts.md section 3): every key term needs a citation key + a quote (<= 15 words)
% copied from the cited paper, same sense; plain / coined otherwise (coined <= 2); "fixed" terms are not exempt.
% concept: <core concept A> | core | borrowed: <citekey> "<quote from the cited paper, same sense>" | <definition>
% concept: <core concept B> | core | composed: <part> (<citekey> "<quote>") + <part> (plain) | <definition>
% concept: <framing term> | framing | borrowed: <citekey> "<quote>" | <definition>
% term: <component / metric / recurring noun> | borrowed: <citekey> "<quote>" | <definition> | other senses: <none or list>
% epithet: <two-word method epithet = negation of the two drawback labels; verbatim in abstract, intro, Methodology opener, Conclusion>
% \showconceptsfalse      % uncomment for the submitted / camera-ready PDF (black text)

% Title: "Verb-ing <Object> with <Key Ingredient>" [most] or
% "<Hook>: <Descriptive Subtitle>" [some]. Sober; no hype, no teasing question.
% Keep this paper's rejected candidates commented here [some].
\title{Verb-ing Object with Key Ingredient}

\author{Anonymous Authors}    % hidden until \iclrfinalcopy

% \iclrfinalcopy % camera-ready only
\begin{document}

\maketitle
\suppressfloats[t]  % [skill choice], not house practice. It keeps the teaser off
                   % page 1 so that it lands at the top of p2, where a one-column
                   % teaser usually settles by float placement alone. Delete it to
                   % allow a page-1 results figure under the abstract [some].

% Logic outline (keep as comments [some]):
%   gap pair -> proposal -> mechanism -> (training signal) -> qualitative result
% Rules: ONE paragraph; NO numbers, NO citations, NO dataset names [all].
% Length at most 250 words [most]; 120-200 at ICML/NeurIPS [current]; 5-9
% sentences (ICML template: 4-6). Expand every acronym; an acronym method name
% gets bold initials [some], a plain name may stay plain [some]. Present tense.
% Code link only at camera-ready.
% Paper-type variants of S6-S7: references/paper_types.md.
\begin{abstract}
% S1 context + difficulty (concessive opener)
Although <task> matters for many real-world applications, <model family> still struggle with it because <reason>.
% S2 current paradigm / named prior work
A recent line of work, <paradigm>, handles <task> by <mechanism>.
% S3 gap as a pair
Yet this <approach> has two key weaknesses: (i) <drawback 1>; and (ii) <drawback 2>.
% S5 proposal
To resolve both, we propose \textbf{X} (\textbf{X}xx \textbf{Y}yy), which <key ingredient in one clause> for more <epithet> <task>.
% S6 mechanism -- pick the variant of the paper type:
%   component method (W): "Specifically, <component A> in X acts only when <condition>, while <component B> <acts> as soon as <condition>."
%   objective method (T): "Specifically, X treats <object> as learnable and trains it under two objectives, <A> and <B>. <A> keeps .... <B> ensures ...."
%   reward method (R):    "Specifically, <the signal is produced by ...>, and <how the signal improves as training proceeds>."
Specifically, X <achieves objective A> through <mechanism A>, and <achieves objective B> through <mechanism B>.
% S7 training signal (training-based methods only [some]):
% We learn these behaviours with <signals> inside <training>, which reward <behavior A> and <behavior B>, respectively.
% S8 qualitative result (firm but bounded)
Experiments with several <models> on <tasks> show that X improves <metric> while preserving <other property>.
% S10 camera-ready only:
% Our code is available at \url{https://anonymous.4open.science/r/XXXX}.
\end{abstract}

% Introduction: 650-750 prose words, 4-5 paragraphs of median 4-6 sentences,
% one teaser, at most one bullet list. NO contribution list, NO "rest of the
% paper" paragraph [all]. No section roadmap here (the intro has none).
% Source conventions: one sentence per line; keep replaced wording as a %
% comment next to the new line (structure.md section 3).
% The slots below are type-neutral. For the P2-P4 wording of your paper type
% (W workflow with components, T objective + constraint with a closed form,
% L loss with theory, R reward/supervision design, F findings paper) open
% references/paper_types.md.
\section{Introduction} \label{sec:introduction}

\begin{figure*}[t]
    \centering
    \figplaceholder[3.4cm]{Teaser (0.4--0.5 pp): the problem / key idea, ideally with evidence panels; concrete example inside}
    \vspace{-8pt}
    \caption{<Title-like phrase>. <What the method does in one sentence>. <The two risks / the takeaway>.}
    \label{fig:teaser-figure}
    \vspace{-12pt}
\end{figure*}

% P1 context -> importance -> current practice -> gap and consequence (2-4 sentences)
<Task> is central to many real-world applications \citep{Placeholder-A, Placeholder-B}, including <application> \citep{Placeholder-C}.
Yet it remains hard in practice: <difficulty 1> \citep{Placeholder-A}, and <difficulty 2> \citep{Placeholder-B}.

% P2 the current paradigm made concrete, with its core formula in the
% Preliminary's symbols and a concrete quoted example [most] -- or the
% challenge + ONE bold-italic question [some]: \textit{\textbf{Can ... ?}} If so, ...
<Paradigm> \citep{Placeholder-A} has recently become a popular way to handle <task> without <cost>.
As Figure~\ref{fig:teaser-figure} illustrates, it takes an input (\eg ``<concrete example>'') and <operation of the prior method> to produce <output>.
In practice, <implementation fact>.

% P3 the PAIR of limitations: bold-headed bullets of 85-115 words each [some],
% or two prevailing strategies + their shared flaw [some], or an enabling
% research line with "On the one hand ... On the other hand" [some]
However, <the prevailing approach> carries over the weaknesses of <root cause> \citep{Placeholder-B}.
In particular, two limitations restrict the <twin concepts> of <task>:
% concepts: mark each key concept once here with \concept{phrase} (+ \citep when borrowed); echo the exact phrase later, unmarked
\begin{itemize}[leftmargin=*]
    \item \textbf{<Limitation 1 as a noun phrase>.} Figure~\ref{fig:teaser-figure} shows that <mechanism>. <Consequence>. Consequently, <gap>.
    \item \textbf{<Limitation 2 as a noun phrase>.} <Mechanism>. For instance, <concrete case (\eg ...)>. The problem grows worse when <condition>.
\end{itemize}

% P4 proposal -> one-sentence key idea with (i)/(ii) -> mechanism with the
% Methodology's inline formula -> (training) -> benefit closer
To this end, we propose a <kind of> approach, \textbf{X}, for <epithet> <task> that borrows the idea of <idea> from <field> \citep{Placeholder-C}.
Our key idea is to <precise verb> <object> with two <components or objectives> --- \concept{<core concept A>} and \concept{<core concept B>} --- which guarantee (i) <property 1> and (ii) <property 2>.
Concretely, <how A and B work, naming each symbol>, written compactly as <core formula in the Methodology's symbols>.
<Optional training sentence, training-based methods only: "We learn these behaviours with <signals> ...">
<A> delivers <benefit A> and removes <limitation 1>, whereas <B> delivers <benefit B>; combined, they make <task> more <epithet>.

% P5 evidence roadmap: one sentence per experiment subsection, same order as RQs
We validate X through extensive experiments.
First, X generally outperforms <baseline> across <tasks> and model sizes (\cf Section~\ref{sec:main-performance}).
Second, we show that <component A> <effect> and that <component B> <effect> (\cf Section~\ref{sec:mechanism}).
Third, an ablation study isolates the contribution of each component (\cf Section~\ref{sec:ablation}).
Overall, X offers a <epithet> recipe for <task> built around <key ingredient>.

% Preliminary: 0.6-0.9 pp (350-550 words), 2 subsections, 2-4 numbered eqs,
% paragraphs of 1-3 sentences between displays. Opens with a roadmap naming
% the subsections [most]; closes with the limitation the method removes [all].
% Base-paradigm displays by paper type (references/paper_types.md):
%   reward design (R):       r = \mathbb{I}(\texttt{match}(\hat{\mathcal{Y}},\mathcal{Y}))
%   operator / closed form (T): \mathbf{a} = \mathbf{W}\mathbf{z}
%   loss design (L):          the base loss \Lapl_{\text{NLL}} written out in full
\section{Preliminary} \label{sec:preliminary}
This section reviews the background our method builds on: <task formulation> in Section~\ref{sec:task-formulation} and <base paradigm> in Section~\ref{sec:base-paradigm}.

\subsection{Task Formulation} \label{sec:task-formulation}
Given an input $\mathcal{X}$ (\eg <concrete input>), the goal is to produce an output $\hat{\mathcal{Y}}$ that matches the ground truth $\mathcal{Y}$.
Formally, <object> is defined as:
\begin{equation} \label{eq:task}
    \hat{\mathcal{Y}} = f_\theta(\mathcal{X}),
\end{equation}
where $f_\theta$ denotes the model parameterized by $\theta$.

\subsection{<Base Paradigm>} \label{sec:base-paradigm}
<Prior work> \citep{Placeholder-A} <does what> by optimizing:
\begin{equation} \label{eq:base-objective}
    \mathcal{J}(\theta) = \mathbb{E}_{(\mathcal{X}, \mathcal{Y}) \sim \mathcal{D}} \big[ s(\hat{\mathcal{Y}}, \mathcal{Y}) \big],
\end{equation}
where $\mathcal{D}$ denotes the training set and $s(\cdot,\cdot)$ is <scoring function>.
This <so-what sentence: what the objective rewards or does>.
% Optional 2.x Literature Review [some] (~0.4-0.5 pp): group prior methods by
% WHICH TERM of the base equation they modify, as bold-headed bullets, then a
% critique paragraph ("Yet most of these designs rest on heuristics ...").
Nevertheless, <limitation 1> and <limitation 2> keep the vanilla approach from reaching the <twin concepts> that <task> needs.

% Methodology: 2.4-2.8 pp (900-1,100 words) at ICLR; paragraphs of median 2-3
% sentences between displays. Roadmap opener [all] (method papers); each
% subsection opens with a purpose clause; the last subsection assembles the
% method. ALL math needed to re-implement the method stays here [all].
% This scaffold is type-neutral. Subsection plans per paper type (W workflow
% with components; T objective + constraint + closed form; L loss + theorem;
% R reward design with roles) are in references/paper_types.md.
\section{Methodology} \label{sec:methodology}
\begin{figure*}[t]
    \centering
    \figplaceholder[3.6cm]{Framework figure (figure*, 0.88--1.0 textwidth), from templates/diagram\_skeleton.tex: white rounded boxes, one filled focal module, LaTeX math labels, compact legend, concrete example}
    \vspace{-8pt}
    \caption{Overview of X. Given <input> $\mathcal{X}$, <step 1>. After that, <step 2>. Finally, <step 3>.}
    \label{fig:main-figure}
    \vspace{-12pt}
\end{figure*}
In this section, we describe \textbf{X} (\textbf{X}xx \textbf{Y}yy), a <epithet> <kind of method> for <goal>.
Section~\ref{sec:method-a} first introduces <the general formulation / component A>.
After that, Section~\ref{sec:method-b} explains how X achieves <objective B>.
Finally, Section~\ref{sec:method-overall} combines these pieces into the complete X.

% Echo the core concepts (same phrase, unmarked) in subsection titles or run-ins (concepts.md).
\subsection{<Core Concept A> with <Technique>} \label{sec:method-a}
To overcome <limitation 1>, we <precise verb> <component A>, following the success of <idea> \citep{Placeholder-C}.
Specifically, <mechanism>, which we formulate as:
\begin{equation} \label{eq:component-a}
    \hat{\mathcal{Y}} = g_{\phi}\big(f_\theta(\mathcal{X})\big),
\end{equation}
where $g_{\phi}$ denotes <component A> parameterized by $\phi$.
This design enables <benefit>.

% OPTIONAL formal results (theory papers only [some]). Numbering: Definition /
% Lemma with plain counters, "Definition 1", "Lemma 1" [current]; a headline
% theorem restatable and numbered within the section, "Theorem 3.1" [some].
% Borrowed results carry the citation in the optional title [some].
% (ICLR: refer to a display as \Eqref{eq:component-a}; NeurIPS/ICML: "Equation~\ref{eq:component-a}".)
We measure <property> with the following quantity:
\begin{definition}[<Named Quantity> \citep{Placeholder-B}] \label{def:concept}
For <objects> $P$ and $Q$ on <space>, the <named quantity> is $D(P \,\|\, Q) \triangleq$ <formula>.
\end{definition}
Solving <problem> directly needs <expensive procedure>, so we rewrite it with the following lemma:
\begin{lemma}[<Descriptive Title> \citep{Placeholder-B}] \label{lemma:simplify}
Let <assumption>. Then <statement>.
\end{lemma}
With this identity, <problem> reduces to <simpler problem>, solved in <cost> (see Appendix~\ref{apdx:proof-lemma} for the proof).
Lemma~\ref{lemma:simplify} then leads to our main result:
\begin{restatable}{thm}{MainTheorem} \label{thm:main}
Let $\Omega_h \triangleq$ <definition> for <parameter> $h > 0$.
Then $\mathcal{J}(\theta)$ \textbf{upper-bounds} <target quantity> up to an additive <term in $h$>.
\end{restatable}
We sketch the argument below and give the complete proof in Appendix~\ref{apdx:proof-main-theorem}.
\begin{proof}
Applying <inequality> to <term> bounds <target quantity> by <expression>. Taking the expectation over <data> gives the claim.
\end{proof}
Theorem~\ref{thm:main} implies that <consequence>. In practice, <knob> is set by <setting> (see Algorithm~\ref{alg:method}).

\subsection{<Core Concept B> via <Technique>} \label{sec:method-b}
With <core concept A> handled by <component A>, we next address <core concept B>.
% One bold run-in paragraph per term, then a "Total ..." closer [some].
% Reward terms are written r^{\text{name}} [some]; losses \Lapl_{\text{Name}} [some].
\textbf{<Term 1>.} <One-sentence purpose>, computed as:
\begin{equation} \label{eq:term-1}
    \ell^{\text{term1}} = \text{<expression>},
\end{equation}
where <symbol> denotes <meaning>. This term rewards <behavior>.

\textbf{<Term 2>.} In addition, to <purpose>, we add:
\begin{equation} \label{eq:term-2}
    \ell^{\text{term2}} = \text{<expression>},
\end{equation}
where <symbol> denotes <meaning>.

\textbf{Total objective.} Finally, we combine the two terms:
\begin{equation} \label{eq:total}
    \ell = \ell^{\text{term1}} + \alpha \, \ell^{\text{term2}},
\end{equation}
where $\alpha$ is a scalar hyperparameter that trades off <trade-off>. The combined objective secures <both properties>.

\subsection{<Method Name / Optimization with X>} \label{sec:method-overall}
Putting the pieces together, X computes the final <output> as:
\begin{equation} \label{eq:final}
    \hat{\mathcal{Y}} = g_{\phi^{\star}}\big(f_{\theta^{\star}}(\mathcal{X})\big).
\end{equation}
Built on <principle> and driven by <signal>, X <achieves A> while <preserving B>.
Algorithm~\ref{alg:method} summarizes the full procedure.
Further implementation details are in Appendix~\ref{apdx:impl-details}.

% Algorithm: main text, single column, [t] for an iterative training loop
% [some]; appendix for a one-shot/closed-form procedure [some].
% RL form of the steps ("sample outputs" / "score them" / "update the
% policy"): references/paper_types.md.
\begin{algorithm}[t]
\captionsetup{labelformat=default,labelsep=colon,name=Algorithm}
   \caption{Training with X}
   \label{alg:method}
\begin{algorithmic}
   \STATE {\bfseries Input:} Dataset $\mathcal{D}$, batch size $B$, and maximum training steps $T$.
   \STATE {\bfseries Initialization:} model $f_\theta$, <component> $g_\phi$, training step $t \leftarrow 0$.
   \WHILE{$t \leq T$}
     \STATE \tikzmk{phaseAa}\CommentBelow{sample a batch and run the model}
     \STATE $\{(\mathcal{X}_n, \mathcal{Y}_n)\}_{n=1}^{B} \sim \mathcal{D}$, \ $\hat{\mathcal{Y}}_n = g_\phi(f_\theta(\mathcal{X}_n))$\tikzmk{phaseAb}
     \STATE \tikzmk{phaseBa}\CommentBelow{compute the objective}
     \STATE $\ell = \ell^{\text{term1}} + \alpha\, \ell^{\text{term2}}$\tikzmk{phaseBb}
     \STATE \tikzmk{phaseCa}\CommentBelow{update parameters}
     \STATE Update $\theta, \phi$ with $\nabla \ell$\tikzmk{phaseCb}
     \STATE $t \leftarrow t + 1$
   \ENDWHILE
   \STATE {\bfseries return} $f_\theta, g_\phi$.
\end{algorithmic}
\boxit[Forward]{PhaseABg}{phaseAa}{phaseAb}
\boxit[Objective]{PhaseBBg}{phaseBa}{phaseBb}
\boxit[Updating]{PhaseCBg}{phaseCa}{phaseCb}
\end{algorithm}

% Experiments: 2.8-3.0 pp at ICLR (850-1,000 words), 1-2 tables + 2-3 figures;
% paragraphs of median 3-4 sentences. The section opens with the RQ list, not a
% roadmap. Opener + RQ bullets -> setup run-ins -> one "(RQk) Title" subsection
% per RQ. No data yet: keep every table fully structured with \pend cells and
% write each finding as a falsifiable expectation inside \draft{} (experiments.md).
\section{Experiments} \label{sec:experiments}
% Opener with an appendix pointer [some]; a bare RQ list without an opener is
% also seen [some].
We evaluate X through the following research questions (further analyses, \eg <x>, <y> and <z>, are in Appendix~\ref{apdx:experiments}):
\begin{itemize}[leftmargin=*]
    \item \textbf{RQ1:} (Performance) Can X achieve better <metric> compared to <baseline> while <keeping cost or property>?
    \item \textbf{RQ2:} (Mechanism) How do <core concept A> and <core concept B> affect the <behavior>?
    \item \textbf{RQ3:} (Ablation Study) How does the hyperparameter $\alpha$ affect X? How does each component contribute to the final performance?
\end{itemize}

% Setup run-ins sit under \subsection{Experiment Setup} [some] or directly
% under the section [some]. Heads: \textbf{Models.} [most], then the data
% paragraph(s): \textbf{Training and evaluation.} [some], \textbf{Evaluation
% Benchmarks.} [some], or one head per benchmark family for an inference-time
% method [some] (e.g. \textbf{In-domain benchmarks.}, \textbf{Transfer benchmarks.}).
\subsection{Experiment Setup} \label{sec:exp-setup}
We summarize the setup here, and Appendix~\ref{apdx:impl-details} gives the full details.

\textbf{Models.} We run all experiments on two backbones of different sizes, <Model-S> and <Model-L> \citep{Placeholder-A}.

\textbf{Training and evaluation.} We train on <dataset> \citep{Placeholder-B}, which holds <size and content>.
For <aspect 1>, we test on <A> \citep{Placeholder-B} and <B> \citep{Placeholder-C}. For <aspect 2>, we test on <C> \citep{Placeholder-C}.
<A> matches the training distribution, and the other test sets are out of distribution.
We report the mean over <n> runs with different random seeds.

% \textbf{Baselines.} only when there are >= 2 baseline families [some],
% grouped by the Related Works taxonomy, plus an ablation baseline. With a
% single natural baseline, delete this paragraph and name the baseline inside
% RQ1 together with the shared starting checkpoint [some; see RQ1 below].
\textbf{Baselines.} We compare against <n $\geq$ 3> <kind> methods from <$\geq$ 2> families: <family 1> methods <Baseline-A> \citep{Placeholder-A} and <Baseline-B> \citep{Placeholder-B}, and the <family 2> method <Baseline-C> \citep{Placeholder-C}. We also include <X w/o C1> as an ablation baseline. Appendix~\ref{apdx:impl-details} describes each baseline.

% Main table: T1 form [some; current]: scale-grouped, no Avg column, one
% metric, our row shaded per cell, no "(Ours)". For many baselines per backbone
% use T2 (Avg column, "+ X (Ours)"); for a quality + cost claim use T4.
\begin{table*}[t]
\centering
\renewcommand{\arraystretch}{0.9}
\caption{Results on <task> benchmarks at two model scales. The best result in each column is \textbf{bold}.}
\makebox[\linewidth][c]{%
  \shrinktofit{0.92\linewidth}{%
    \begin{tabular}{cl|cccc}
    \toprule
    \multicolumn{2}{c|}{\textbf{}} & \multicolumn{4}{c}{\textbf{Benchmarks}} \\
    \textbf{Scale} & \textbf{Method} & \textbf{Bench-1} & \textbf{Bench-2} & \textbf{Bench-3} & \textbf{Bench-4} \\
    \midrule
    \multirow{4}{*}{Model-S}
    & Base       & \pend & \pend & \pend & \pend \\
    \cmidrule(lr){2-6}
    & + Baseline-A & \pend & \pend & \pend & \pend \\
    & + Baseline-B & \pend & \pend & \pend & \pend \\
    & \cellcolor{oursrow}+ X & \cellcolor{oursrow}\pend & \cellcolor{oursrow}\pend & \cellcolor{oursrow}\pend & \cellcolor{oursrow}\pend \\
    \midrule\midrule
    \multirow{4}{*}{Model-L}
    & Base       & \pend & \pend & \pend & \pend \\
    \cmidrule(lr){2-6}
    & + Baseline-A & \pend & \pend & \pend & \pend \\
    & + Baseline-B & \pend & \pend & \pend & \pend \\
    & \cellcolor{oursrow}+ X & \cellcolor{oursrow}\pend & \cellcolor{oursrow}\pend & \cellcolor{oursrow}\pend & \cellcolor{oursrow}\pend \\
    \bottomrule
    \end{tabular}}}
\label{tab:main-table}
\vspace{-16pt}
\end{table*}

\subsection{(RQ1) Performance Comparison} \label{sec:main-performance}
Table~\ref{tab:main-table} reports the results of X and the baselines, and Appendix~\ref{apdx:experiments} adds results under <other setting>.
% Single-baseline variant [some]: "Table~\ref{tab:main-table} compares X with <baseline> trained on <signal> alone; both runs start from one shared <SFT> checkpoint."
We have the following observations:
\begin{itemize}[leftmargin=*]
    \item \textbf{\draft{[Expected]} X generally outperforms <baseline> on every benchmark family.}
    As shown in Table~\ref{tab:main-table}, <evidence>. Moreover, <strengthening>.
    We attribute this gain to <mechanism>.
    In contrast, <baselines> are weaker and less stable, probably because <reason>.
    \item \textbf{\draft{[Expected]} X is more <efficient> than <baseline>.}
    <Evidence with a bounded number, \eg ``up to''>.
\end{itemize}

\subsection{(RQ2) Mechanism of <Components>} \label{sec:mechanism}
In this section, we examine how <components> shape <behavior>.

\begin{figure}[t]
    \centering
    \begin{subfigure}[t]{0.49\linewidth}
        \centering
        \figplaceholder[2.6cm]{Plot: <metric> dynamics (one PNG per panel)}
        \vspace{-6pt}
        \caption{<Model-S>.}
        \label{fig:mechanism-a}
    \end{subfigure}
    \hfill
    \begin{subfigure}[t]{0.49\linewidth}
        \centering
        \figplaceholder[2.6cm]{Plot: <metric> dynamics (one PNG per panel)}
        \vspace{-6pt}
        \caption{<Model-L>.}
        \label{fig:mechanism-b}
    \end{subfigure}
    \vspace{-8pt}
    \caption{<Title phrase>. <Takeaway sentence>.}
    \label{fig:mechanism}
    \vspace{-12pt}
\end{figure}

\textbf{<Core concept A>.} To test whether <core concept A> lowers <risk>, we monitor <quantity>.
As shown in Figure~\ref{fig:mechanism}, <observation>. This phenomenon is because <mechanism>. Consequently, <consequence>.

\textbf{<Core concept B>.} We build <stress setting> to isolate <core concept B>. <Observation>. More results are in Appendix~\ref{apdx:experiments}.

\subsection{(RQ3) Ablation Study} \label{sec:ablation}
We remove <component> and evaluate the resulting variant on <Model-L> (\ie X (w/o <component>)).
\textbf{\draft{[Expected]} Both components contribute to the final performance.} <Evidence>.
We therefore set $\alpha$ to <value> by default, because <reason>. Appendix~\ref{apdx:experiments} reports further results.

% % Kept empty on purpose: Related Works lives in Appendix A [most]; see
% chapters/7_appendix.tex. A findings/analysis paper may keep it in the main
% text as Section 2 [some]; then move Appendix A here unchanged.
   % Related Works lives in Appendix A [most]
% Limitations: numbered section BEFORE the Conclusion, 2-3 sentences, 35-60
% words, two items, no mitigation, no promise [most]. Items: preferably
% "On the one hand ... On the other hand ..." [current], otherwise
% "First ... In addition ..." [older]. Present tense.
\section{Limitations} \label{sec:limitations}
This work has two main limitations.
On the one hand, <extra cost or stability issue>.
On the other hand, <scope limit>, so whether X transfers to <other domain> remains untested.

% Conclusion: ONE paragraph, 80-140 words, mirrors the abstract, no new claims;
% 3-6 sentences; past tense [most] or present [some] -- one tense per paragraph.
\section{Conclusion} \label{sec:conclusion}
Prior <paradigm> methods such as <Prior> \citep{Placeholder-A} bring <benefit>, but they suffer from <problem 1> caused by <cause> and from <problem 2> because they lack <mechanism>.
We proposed X, which <achieves objective A> through <mechanism A> and <achieves objective B> through <mechanism B> (\ie <core concept A> and <core concept B>).
Experiments on several models confirmed that X is effective and makes <task> more <epithet>.
Future work will extend X to <domain>.
% Optional footnote on the last sentence pointing to a statement (2026-style ICLR kits:
% "LLM use is described in Appendix~\ref{apdx:use-of-llms}"; NeurIPS: broader impacts). ICLR 2027
% discloses LLM use in the AI Use Statement before the references instead.

\label{main-text-end}  % page check: this label's page must equal the page limit (see SKILL.md)

% Camera-ready only (ICLR kit: unnumbered third-level heading; NeurIPS: \begin{ack}):
% \subsubsection*{Acknowledgments}
% This research is supported by <grant>.
% \clearpage
      % Limitations + Conclusion (+ Acknowledgments)
% ICLR only: unnumbered statements after Conclusion/Acknowledgments and before
% the references, not counted (ICLR 2027 kit l.397-444). NeurIPS: Broader
% Impacts in the appendix + checklist. ICML: \clearpage \section*{Impact
% Statement} (main_skeleton_icml.tex). Details: references/venues.md §1.5.
%
% ICLR 2027: REQUIRED, <= 1 page (K-ICLR27 l.397-412; iclr.cc 2027 AI Policy).
% Disclosure is REQUIRED if generative AI was used to: generate synthetic data;
% develop theoretical models or conceptual frameworks; formulate mathematical
% claims; provide critical ingredients for proofs or help write them; propose or
% refine hypotheses; design or give feedback on the research methodology or
% experiments; implement methods; assist with translation; clean or reformat
% data; support qualitative and thematic data analysis; interpret results.
% RECOMMENDED for: drafting or editing text, figures, code edits, experimental
% parameters, literature search and summaries, brainstorming, references, titles.
% The kit's text has four parts: required tasks AI WAS used for, required tasks
% it was NOT used for, required tasks that do not apply, then "We have reviewed
% all AI-assisted work" + how it was checked + responsibility. A 2026-style
% statement that names only ideation, experiment design, data analysis and
% interpretation misses implementation, synthetic data, proofs, translation and
% data cleaning: name every required category below.
\section*{AI Use Statement}
We did not use generative AI tools to generate data, develop theory, formulate or prove mathematical claims, propose or refine hypotheses, design or give feedback on the methodology or experiments, implement methods, translate, clean or reformat data, analyse data, or interpret results. We used them to <copy-edit the wording of Sections 1-3> and <draft plotting scripts>. We have reviewed all AI-assisted work (<how, e.g., all plotting code was run and its output checked by one author>) and take responsibility for the final content of this paper.

% Ethics: three short paragraphs [some] -- aim and data, possible misuse, limits of the guarantee.
\section*{Ethics Statement}
This work aims to <goal>, which can benefit <stakeholders>. Our experiments use only public models and datasets, and they involve neither private data nor human subjects.

We also recognize a possible misuse: <technique> might, for instance, be used to <misuse>.

X improves <property>, but it cannot fully guarantee <guarantee>. We encourage further work on <direction>.

% Reproducibility: ONE paragraph that points to sections, appendix and code; it does not repeat them.
\section*{Reproducibility Statement}
We take several steps to make our results reproducible. An anonymous implementation is available at \url{https://anonymous.4open.science/r/XXXX}. Section~\ref{sec:methodology} describes the method and Section~\ref{sec:exp-setup} the experimental setup. Appendix~\ref{apdx:impl-details} lists the implementation details, hardware, hyperparameters, data selection and prompts. Additional analyses appear in Section~\ref{sec:experiments} and Appendix~\ref{apdx:experiments}.
   % ICLR only; before references

\bibliography{references}
\bibliographystyle{iclr2027_conference}

% Appendix: every top-level section starts with \clearpage and uses an apdx:
% label. Section order: references/structure.md (appendix order).
\appendix
\clearpage
\section{Related Works} \label{apdx:related-works}
% ~1 page, 2-3 subsections mapped onto the method's ingredients; each
% paragraph: significance -> taxonomy -> representatives -> shared
% limitation -> bridge. The last subsection ends with the positioning sentence.
\subsection{<Base Paradigm>}
<Field> is now a standard tool for <goal> \citep{Placeholder-A}. Most existing methods fall into two groups, <A> \citep{Placeholder-B} and <B> \citep{Placeholder-C}. Methods of the first group ... Methods of the second group ... Although both groups work well in their settings, they are prone to <problem 1> and <problem 2>.

\subsection{<Enabling Technique>}
<Technique> has recently shown promise for <goal> \citep{Placeholder-A, Placeholder-B}. ... This progress motivates our design of <component>.

\subsection{<Form of Supervision / Constraint>}
... Yet how to <open problem> is still an open question. We close this gap with a framework that <does>, combining <A> with <B>.

\clearpage
\section{Theoretical Analysis} \label{apdx:theory}
% Only when the main text has a formal result [some].
\subsection{Proof of Lemma~\ref{lemma:simplify}} \label{apdx:proof-lemma}
We first rewrite <problem> as <equivalent form> using <identity>. Under <assumption>, <step 1>. Substituting <expression> then gives <step 2>, which is <statement>.

\subsection{Proof of Theorem~\ref{thm:main}} \label{apdx:proof-main-theorem}
\MainTheorem*
\begin{proof}
\textbf{Setup.} Let <symbol> denote <object>.

\textbf{Rewriting the objective.} Under the constraint <condition>, <quantity> becomes:
\begin{align}
    a & = b \notag \\
      & = c,
\end{align}
where the second step follows from <justification>.

\textbf{Bounding the target.} Because <fact>, applying <inequality> gives:
\begin{equation}
    <\text{target quantity}> \;\le\; \mathcal{J}(\theta) + C\,h,
\end{equation}
where $C$ depends only on <constant>. Hence <consequence>.
\end{proof}

\clearpage
\section{Implementation Details} \label{apdx:impl-details}
All experiments use the <framework>\footnote{\url{https://github.com/...}} framework on <hardware>.

\textbf{Training.} <Recipe in prose with the key values: learning rate = <lr>, batch size = <batch size>, ...>. <Why a non-obvious setting was chosen>.

\begin{table}[h]
\centering
\caption{Key hyperparameters used to train X.}
\begin{tabularx}{0.45\linewidth}{l|>{\centering\arraybackslash}X}
\toprule
\multicolumn{1}{c}{\textbf{Hyperparameter}} \vline & \multicolumn{1}{c}{\textbf{Value}} \\
\midrule
Learning rate          & \pend \\
Batch size             & \pend \\
Epochs                 & \pend \\
Max sequence length    & \pend \\
Random seeds           & \pend \\
\bottomrule
\end{tabularx}
\label{tab:hyper-train}
\end{table}

\clearpage
\section{Prompts} \label{apdx:prompts}
\begin{figure}[ht]
\centering
\begin{lrbox}{\PromptBoxA}
\begin{gptpromptbox}
\begin{minipage}{0.95\linewidth}\raggedright\scriptsize\ttfamily
You will rate one training example for difficulty. ... \\
\textcolor{roleB}{<score>} ... \textcolor{roleB}{</score>}
\end{minipage}
\end{gptpromptbox}
\end{lrbox}
\usebox{\PromptBoxA}
\caption{Prompt used for the <role> role.}
\label{fig:prompt-role}
\end{figure}

\clearpage
\section{Additional Experiments} \label{apdx:experiments}
% Subsections mirror the main-text RQs: 1-3 sentences of setup, one table or
% figure, 1-3 sentences of conclusion. Plain headers, generic captions.
\subsection{Results on Other Model Sizes}
We also apply X to <models>. The gains hold at every size we test.

\clearpage
\section{Case Study} \label{apdx:case-study}
\begin{figure}[ht]
\centering
\begin{promptbox}
\small
\textbf{Question:} \quad <question> \\
\textbf{$w/o$ X:} \quad \textcolor{bad}{<failure>} \\
\textbf{$\phantom{/o}w$ X:} \quad \textcolor{good}{<success>}
\end{promptbox}
\caption{A case study of X on <task>.}
\label{fig:case-study}
\end{figure}

% ICLR 2026-style kits only: "Use of LLMs" as the last appendix section.
% ICLR 2027 replaced it with a required AI Use Statement BEFORE the references
% (chapters/8_ethics_reproducibility_statement.tex; references/venues.md §1.5).
% \clearpage
% \section{Use of LLMs} \label{apdx:use-of-llms}
% % 2026-style ICLR kits: always the last appendix section.
% LLMs played no part in forming research ideas, designing experiments, or analysing data. We used them only for:
% \begin{itemize}[leftmargin=*]
%     \item <light editing of a few sentences>;
%     \item <help with plotting scripts>.
% \end{itemize}
% Every result, analysis and conclusion in this paper was produced and checked by the authors, who are fully responsible for its content.

% NeurIPS: \clearpage \input{checklist}

\end{document}
